% Part: incompleteness
% Chapter: arithmetization-syntax
% Section: coding-terms

\documentclass[../../../include/open-logic-section]{subfiles}

\begin{document}

\olfileid[id]{inc}{art}{trm}
\olsection{Pengodean Suku}

\begin{explain}
Suku hanyalah suatu jenis barisan simbol tertentu: suku dibangun secara
induktif dari konstanta dan variabel menurut aturan pembentukan suku. Karena
barisan simbol dapat dikodekan sebagai bilangan---dengan menggunakan skema
pengodean simbol serta suatu cara untuk mengodekan barisan bilangan---tidak
sulit untuk memasangkan bilangan G\"odel dengan suku. Tantangannya justru
menunjukkan bahwa sifat yang dimiliki suatu bilangan apabila bilangan itu
merupakan bilangan G\"odel dari suku yang terbentuk dengan benar bersifat dapat
dihitung, atau bahkan rekursif primitif.
\end{explain}

!!^{variable}s dan !!{constant}s merupakan suku yang paling sederhana, dan
mudah untuk menguji apakah $x$ adalah bilangan G\"odel dari suku semacam itu:
$\fn{Var}(x)$ berlaku jika $x$ adalah $\Gn{\Obj v_i}$ untuk suatu~$i$. Dengan
kata lain, $x$ adalah barisan dengan panjang~$1$ dan satu-satunya elemennya,
$(x)_0$, merupakan kode suatu !!{variable}~$\Obj v_i$, yakni $x$ adalah
$\tuple{\tuple{1, i}}$ untuk suatu~$i$. Demikian pula, $\fn{Const}(x)$ berlaku
jika $x$ adalah $\Gn{\Obj c_i}$ untuk suatu~$i$. Kedua relasi ini rekursif
primitif, sebab jika $i$ semacam itu ada, maka indeks tersebut harus $< x$:
\begin{align*}
  \fn{Var}(x) & \defiff \bexists{i<x}{x = \tuple{\tuple{1, i}}}\\
  \fn{Const}(x) & \defiff \bexists{i<x}{x = \tuple{\tuple{2, i}}}
\end{align*}

\begin{prop}
\ollabel{prop:term-primrec}
Relasi $\fn{Term}(x)$ dan $\fn{ClTerm}(x)$, yang masing-masing berlaku jika dan
hanya jika $x$ adalah bilangan G\"odel suatu suku atau suku tertutup, bersifat
rekursif primitif.
\end{prop}

\begin{proof}
Suatu barisan simbol~$s$ adalah suku jika dan hanya jika terdapat barisan suku
$s_0$, \dots, $s_{k-1} = s$ yang mencatat bagaimana suku~$s$ dibentuk dari
!!{constant}s dan !!{variable}s menurut aturan pembentukan suku. Agar barisan
yang disangka sebagai barisan pembentukan itu mengikuti aturan pembentukan,
untuk setiap $i < k$ harus berlaku salah satu dari hal berikut:
\begin{enumerate}
\item $s_i$ adalah !!a{variable}~$\Obj v_j$; atau
\item $s_i$ adalah !!a{constant}~$\Obj c_j$; atau
\item $s_i$ dibangun dari $n$ suku $t_1$, \dots, $t_n$ yang muncul sebelum
  posisi~$i$, dengan menggunakan !!{function} beraritas $n$,~$\Obj f^n_j$.
\end{enumerate}
Untuk menunjukkan bahwa relasi yang bersesuaian pada bilangan G\"odel bersifat
rekursif primitif, kita harus menyatakan syarat ini secara rekursif primitif,
yakni dengan menggunakan fungsi dan relasi rekursif primitif serta kuantifikasi
terbatas.

Andaikan $y$ adalah bilangan yang mengodekan barisan $s_0$, \dots,
$s_{k-1}$, yakni $y = \tuple{\Gn{s_0}, \dots, \Gn{s_{k-1}}}$. Bilangan itu
mengodekan barisan pembentukan bagi suku dengan bilangan G\"odel~$x$ jika dan
hanya jika untuk setiap $i < k$:
\begin{enumerate}
\item $\fn{Var}((y)_i)$; atau
\item $\fn{Const}((y)_i)$; atau
\item terdapat $n$ dan bilangan~$z = \tuple{z_1, \dots, z_n}$ sedemikian rupa
  sehingga setiap $z_l$ sama dengan suatu $(y)_{i'}$ untuk $i' < i$, dan
\[
(y)_i = \Gn{\Obj f^n_j(} \concat \fn{flatten}(z) \concat \Gn{)},
\]
\end{enumerate}
serta $(y)_{k-1} = x$. (Fungsi $\fn{flatten}(z)$ mengubah barisan
$\tuple{\Gn{t_1}, \dots, \Gn{t_n}}$ menjadi $\Gn{t_1, \dots, t_n}$ dan
bersifat rekursif primitif.)

Indeks $j$ dan $n$, bilangan G\"odel $z_l$ dari suku $t_l$, serta kode~$z$
dari barisan~$\tuple{z_1, \dots, z_n}$ dalam (3), semuanya lebih kecil
daripada~$y$. Kita dapat mengganti $k$ di atas dengan $\len{y}$. Dengan
demikian, kita dapat menyatakan ``$y$ adalah kode suatu barisan pembentukan
bagi suku dengan bilangan G\"odel~$x$'' dengan cara yang menunjukkan bahwa
relasi ini rekursif primitif.

Sekarang kita hanya perlu meyakinkan diri bahwa terdapat batas rekursif
primitif bagi~$y$. Jika $x$ merupakan bilangan G\"odel suatu suku, suku itu
harus memiliki barisan pembentukan yang memuat paling banyak $\len{x}$ suku
(karena setiap suku dalam barisan pembentukan~$s$ harus dimulai pada suatu
posisi dalam~$s$, dan tidak ada dua sub-suku yang dapat dimulai pada posisi
yang sama). Tentu saja bilangan G\"odel setiap sub-suku~$s$ ialah $\le x$.
Jadi, selalu terdapat barisan pembentukan dengan kode
$\le p_{k-1}^{k(x+1)}$, dengan $k=\len{x}$.

Untuk $\fn{ClTerm}$, cukup hilangkan klausa bagi !!{variable}s.
\end{proof}

\begin{prob}
Tunjukkan bahwa fungsi $\fn{flatten}(z)$, yang mengubah barisan
$\tuple{\Gn{t_1}, \dots, \Gn{t_n}}$ menjadi $\Gn{t_1, \dots, t_n}$,
bersifat rekursif primitif.
\end{prob}

\begin{prop}\ollabel{prop:num-primrec}
  Fungsi $\fn{num}(n) = \Gn{\num{n}}$ bersifat rekursif primitif.
\end{prop}

\begin{proof}
  Kita mendefinisikan $\fn{num}(n)$ melalui rekursi primitif:
  \begin{align*}
    \fn{num}(0) & = \Gn{\Obj 0}\\
    \fn{num}(n+1) & = \Gn{\prime(} \concat \fn{num}(n) \concat \Gn{)}.
  \end{align*}
\end{proof}

\end{document}
